\section{Conclusion}
\label{sec:conclusion}

The great promise of \aisub{s} lies in their potential to accurately predict human behavior in novel settings.
Realizing this capability could transform social science research and public policy.
It could provide the social science equivalent of a lab bench in the physical sciences: an accurate, scalable playground to test ideas before large-scale and expensive implementation with humans.\footnote{This could be even more powerful given the often-observed researcher inability to predict results of their own experiments \citep{predictscience2019DellaVigna,milkmanMegastudies2021,Gandhi2023Hypothetical,Gandhi2024EffectSize,Duckworth2025ZearnMega}.}
Yet, with the current state-of-the-art foundation models, \aisub{s} are not yet reliable enough to be used in this way out of the box.

In this paper, we explored an approach to address this shortcoming.
Our approach relies on two key principles: (i) grounding candidate \aisub{s} in theories expected to drive human behavior in the target setting, and (ii) optimizing and then validating \aisub{s} in distinct but related settings presumed to be well-explained by the same theory.
Without theoretical grounding, optimized \persona{s} may fail to meaningfully improve even in-sample predictions.
Without validation across distinct but related datasets, optimized \persona{s} are prone to overfit a particular data-generating process.
Just as economists carefully extend established theories to novel policy contexts---relying on accumulated empirical validation rather than absolute certainty---optimizing theoretically-grounded \aisub{s} to match samples of human data, and then validating them in distinct but related settings, provides a principled foundation for predicting humans in new settings.

The improvements in predictive power yielded by this approach are substantial. 
In four novel and preregistered strategic games derived from \citeauthor{1120Arad2012}'s 11-20 money request game, theory-grounded \aisub{s}, optimized and validated through our methodology, reduced prediction errors by approximately \MinPercImproveArad-\MaxPercImproveArad\% compared to baseline AI predictions. 
Remarkably, these theory-grounded agents predicted initial play in some of the games better than the original human data from \citeauthor{1120Arad2012}. 
Importantly, our results are not confined to games involving strategic reasoning.
In Appendix \ref{app:predict_new_CR}, we apply the same procedure to the allocation games from \citeauthor{charness2002understanding}, and, using data from a preregistered experiment with entirely novel human responses, we find substantial improvements.

Although this approach provides no statistical guarantees for arbitrary novel settings---indeed, no procedure can guarantee predictive power in entirely novel domains without a fully specified and correct causal model---we demonstrated that we can make externally valid inferences within a pre-committed family of settings.
Using novel data from \SampleSizeHuman{} participants playing \SampleSizeGame{} games randomly sampled from a heterogeneous population of \TotalGames{} strategic games, we found that the strategic level-$k$ agents generalized effectively across this broad domain.
In \OptGameAllMatchRate\% of games, all human subjects chose actions within the support of the optimized \aisub{} responses. 
These general agents substantially outperformed the baseline AI off-the-shelf, a cognitive hierarchy model, and the Harsanyi-Selten equilibria.

The results were also robust to several alternative specifications and were consistent across different game structures.
Because the games were randomly sampled under the assumptions stated in Section~\ref{sec:estimation}, these results are externally valid for the entire population of \TotalGames{} games.
While we can only guarantee validity within this specific population and sampling distribution, the strong performance of theoretically motivated \aisub{s} across such a diverse set of strategic games suggests that similar approaches would likely generalize to alternative distributions and even other strategic games.

Looking ahead, researchers could leverage this approach by identifying broad domains where they need predictions and can obtain training and validation data from representative subsets---even if there is substantial structural variation within the domain.
Indeed, a key advantage of \aisub{s} is that they can make predictions in response to any setting expressed in natural language without any prior data.
This is often not possible with traditional machine learning methods and economic models.

Another exciting research direction would be to automate the theory-prediction-testing loop.
Such a system would start with novel human data and relevant experimental settings, iteratively generate theory-informed candidate personas, optimize their parameters, and systematically evaluate their generalizability across samples.
One can also imagine just starting with a novel setting, and then having an AI system actually search for the relevant training and validation data to execute our approach.
Recent research supports the feasibility of these ideas \citep{Jackson2025Mixture,zhu2025ExplPredicting}.
In particular, \cite{Manning2024Automated} demonstrate AI systems capable of automating the full social scientific workflow---from hypothesis generation to experimental simulation and data analysis.

More generally, the results in this paper linking theory-grounded \aisub{s} to robust generalizability in novel settings, and atheoretical \aisub{s} to failure, are notable for reasons beyond prediction \citep{hofman2021integrating}.
They suggest that the underlying LLM has correctly learned the relevant relationships between the \aisub{s} and human responses to the given setting.
This is even more notable given that it is highly unlikely that these mappings were explicitly specified during training.  
Such a finding aligns with recent evidence that LLMs form rich internal representations of their human-generated training corpus rather than merely memorizing it \citep{lindsey2025biology, ameisen2025circuit}.
If true for LLMs and human behavior more generally, our \persona{}-alignment-generalizability exercises may offer more than improved predictive power.
If an LLM armed with a particular theory-grounded \persona{} matches human data particularly well across a wide range of related settings, it might be evidence that the theory has a lot of explanatory power for the underlying human sample.
Building on a young social scientific literature \citep{peterson2021lottery,hirasawa2022using,EnkeComplexity2023,Si2024NovelResearch,MLhypothesis2024,Ashesh2024Anomaly,batista2024words,movva2025sparse}, this could, in turn, provide researchers with robust machine learning methods to rapidly and efficiently inform promising new hypotheses.

\newpage \clearpage

\bibliographystyle{ecta}
\bibliography{optimize}

\newpage \clearpage

\appendix

\renewcommand{\thefigure}{A\arabic{figure}} 
\setcounter{figure}{0}  

\renewcommand{\thetable}{A\arabic{table}} 
\setcounter{table}{0}  

